% !TEX root = ../main.tex

\chapter*{\MakeUppercase{RESEARCH ON DYNAMIC KAFKA TOPIC DISCOVERY AND SUBSCRIPTION MECHANISM BASED ON FLINK}}
\addcontentsline{toc}{chapter}{英文大摘要}
\markboth{英文大摘要}{英文大摘要}

With the continuous development of real-time data infrastructure, Apache Kafka and Apache Flink have become a common technical combination for event-driven systems, log processing, monitoring, and real-time analytics. Kafka provides durable and high-throughput distributed logs, while Flink provides stateful stream processing, checkpoint-based fault tolerance, and flexible external output through the unified Sink API. In production systems, upstream services write events into Kafka, Flink jobs consume and transform these streams, and the results are written back to Kafka topics or other downstream systems. This architecture works well when the topology is stable, but it becomes less flexible when downstream topics change frequently. In multi-tenant or evolving business systems, Kafka topics may be created by tenant, date, region, business domain, or migration stage. A long-running Flink job may therefore face a changing set of write targets.

The conventional Flink Kafka Sink is mainly configured statically. Users usually specify the target topic, bootstrap servers, serialization schema, and delivery guarantee before the job is submitted. After deployment, these write targets normally remain fixed. If the topic set changes, operators often need to stop the job, modify configuration, resubmit the job, and wait for recovery. This increases operational cost and complicates consistency management. The source side has explored dynamic topic subscription through pattern-based discovery, but the sink side has additional challenges. It must discover target topics, bind records to physical topics, create and reuse write channels, snapshot route state, and coordinate commit semantics.

This thesis studies a dynamic Kafka topic discovery and subscription mechanism based on Flink, with a focus on the Sink side. The goal is to allow a Flink job to discover matching Kafka topics and adjust its write paths without being restarted. The implementation is built as a connector prototype whose core capabilities cover metadata discovery, runtime routing, write-channel management, route-level state snapshotting, and commit coordination. The validation jobs generate test records, distribute route updates, and compare dynamic writing with a static baseline.

The main design idea is the separation between logical routes and physical Kafka topics. Instead of forcing business records to carry a fixed topic name, the system allows records to carry a logical route identifier. This identifier is mapped to a target description consisting of a Kafka cluster, connection information, and a topic pattern. The runtime routing layer resolves the logical destination into one or more physical topics by querying Kafka metadata and matching topic names with a regular expression. This design allows the physical target of a business route to change while keeping upstream business logic stable.

The overall architecture can be described as a dynamic shell over a static core. The static core is the existing Kafka writing capability: after a physical route has been resolved, serialization, producer sending, flushing, and transaction handling are delegated to the mature Kafka Sink path. The dynamic shell is responsible for metadata refresh, route resolution, write-channel creation, route cache management, state snapshotting, and commit preparation. This design avoids rewriting a complete Kafka producer layer and concentrates dynamic behavior in the route management layer.

The metadata discovery layer periodically obtains the currently visible topic set from Kafka's management plane. The discovery result is filtered by a configured regular expression, and the refresh interval bounds the observable discovery delay together with metadata request time and Flink task scheduling delay. The runtime then maintains two write paths. The automatic discovery path writes ordinary records to all currently active routes. The explicit logical route path resolves a record's route identifier into physical topics and writes only to the corresponding target set. Route update events can add, overwrite, or remove route mappings, and the resolved cache is invalidated accordingly.

State and commit management are essential for compatibility with Flink. The dynamic sink stores route-level recovery units instead of a single global state. Each recovery unit binds route metadata, producer configuration, and the underlying write state so that known write channels can be rebuilt during recovery and then aligned with the current Kafka topology. For commit management, the system wraps lower-level commit results with route metadata, groups commit requests by route, and delegates them to the corresponding Kafka commit path. Thus, the transaction path is organized by route-level states and commits rather than one global transaction across all routes.

The thesis defines measurable requirements before presenting experiments. Functionally, the system should support dynamic topic discovery, automatic subscription of writable routes, dynamic routing through logical route identifiers, and stable processing with serialization, Kafka writing, checkpointing, and recovery. Non-functionally, the prototype should satisfy concrete criteria in a local validation environment. With a two-second discovery interval, the first write to a newly created matching topic should occur within three seconds. In the explicit two-route scenario, throughput loss compared with the static baseline should be less than five percent. When route count grows from two to four, throughput should not drop by more than ten percent. With five-second checkpoint intervals, steady checkpoint completion time should stay below one hundred milliseconds. For the transaction path, a small-scale exactly-once run should complete at least five checkpoints and produce committed output on both target topics.

The discovery timeliness experiment generates records without a logical route identifier, starts with a seed topic, and creates new matching topics while the job is running. By polling Kafka offsets, the experiment measures the delay between topic creation and the first positive offset. Six samples are collected in two rounds. The observed first-write delays range from about 1.074 seconds to 2.418 seconds, with an average of 1.530 seconds. These results are consistent with the configured two-second discovery interval and satisfy the three-second target.

The explicit route and throughput experiments load two logical routes from YAML and map them to different topic patterns. In a thirty-second local run with a five-millisecond emission interval, the dynamic two-route job writes 5742 records. The two target topics receive 2812 and 2930 records, corresponding to 49.0 percent and 51.0 percent of the dynamic output. The static baseline writes 5776 records, or about 192.5 records per second. The dynamic two-route job reaches about 191.4 records per second, a decrease of about 0.59 percent, which is below the five-percent threshold. A four-route dynamic run achieves about 192.1 records per second, showing no observable throughput degradation in the small-scale setting.

Checkpoint coordination is validated with five-second checkpoint intervals. Five successful checkpoints are observed, with completion times of 476 ms, 11 ms, 22 ms, 19 ms, and 21 ms. Excluding the cold-start checkpoint, the steady average is about 18.2 ms, below the one-hundred-millisecond threshold. The transaction path is further validated in exactly-once mode. With a unique transactional id prefix and a local transaction timeout setting, five checkpoints complete successfully, no commit failure is observed, and two target topics both receive committed output. This shows that the route-level commit path is executable in a small local environment.

The implementation and experiments demonstrate that extending Kafka Sink with dynamic topic discovery, logical route mapping, route-level write-channel management, and route-level state and commit wrapping is feasible. Compared with rewriting a full Kafka producer layer, this approach reuses existing Kafka Sink behavior and keeps dynamic logic in the connector shell. The prototype can discover new topics, route records according to logical identifiers, write to multiple topics, participate in checkpointing, and execute a small-scale transaction path. It also has limitations. Metadata discovery relies on periodic polling, the automatic discovery path may cause write amplification, retired write channels are not automatically reclaimed in exactly-once mode, and the transaction prefix strategy still deserves stronger collision-resistant design. Future work should focus on metadata efficiency, transaction-safe cleanup, observability, fault injection, and larger distributed validation.
